%==============================================================================
%  CHAPTER 5 — VALIDATION AND RESULTS
%==============================================================================
\chapter{Validation and Results}
\label{chap:validation}

\section{Introduction}
This chapter validates the platform in two complementary ways. First, a
functional validation: every capability delivered by the sprints was exercised
against acceptance scenarios. Second, an end-to-end validation on a
\emph{real} use case: a telecom churn model, trained, tracked, registered,
deployed, and consumed entirely through the platform.

% ---------------------------------------------------------------------------
\section{Functional Validation}
\label{sec:functional-validation}
Each increment was validated manually against its acceptance scenarios before
its sprint closed. Tables~\ref{tab:tests-s1}--\ref{tab:tests-s4} consolidate the
main scenarios, one table per sprint.

\begin{table}[H]\centering
\renewcommand{\arraystretch}{1.15}
\caption{Sprint 1 test scenarios --- Authentication and projects}\label{tab:tests-s1}
\small
\begin{tabularx}{\textwidth}{|X|X|c|}
\hline
\rowcolor{primary!15}
\textbf{Scenario} & \textbf{Expected result} & \textbf{Status}\\
\hline
Sign up, then sign in & Tokens issued; redirect to dashboard & OK\\ \hline
Sign in with wrong password & \texttt{401}, no token issued & OK\\ \hline
Access token expiry & Session silently renewed via refresh token & OK\\ \hline
Create a project & Stored and bound to a new tracking experiment & OK\\ \hline
Open another user's project & \texttt{404}, resource invisible & OK\\ \hline
\end{tabularx}
\end{table}

\begin{table}[H]\centering
\renewcommand{\arraystretch}{1.15}
\caption{Sprint 2 test scenarios --- Training and tracking}\label{tab:tests-s2}
\small
\begin{tabularx}{\textwidth}{|X|X|c|}
\hline
\rowcolor{primary!15}
\textbf{Scenario} & \textbf{Expected result} & \textbf{Status}\\
\hline
Upload a \texttt{.zip} without an entry script & Entry script auto-detected by priority order & OK\\ \hline
Run a scikit-learn script & Metrics, parameters, and model captured automatically & OK\\ \hline
Run a real-world Colab notebook & Converted at upload, shims applied, run succeeds & OK\\ \hline
Run a script with a missing dependency & Dependency installed in the pod before execution & OK\\ \hline
Run a failing script & Run marked failed; readable reason shown on the run & OK\\ \hline
Watch logs during a run & Log lines stream live to the browser & OK\\ \hline
Open a run's artifacts & Folders navigable; images preview; files download & OK\\ \hline
Request another user's artifacts & \texttt{404}, the run is not revealed & OK\\ \hline
\end{tabularx}
\end{table}

\begin{table}[H]\centering
\renewcommand{\arraystretch}{1.15}
\caption{Sprint 3 test scenarios --- Registry and deployment}\label{tab:tests-s3}
\small
\begin{tabularx}{\textwidth}{|X|X|c|}
\hline
\rowcolor{primary!15}
\textbf{Scenario} & \textbf{Expected result} & \textbf{Status}\\
\hline
Register a run, promote to Production & Version and stage updated in registry and interface & OK\\ \hline
Deploy a registered model & InferenceService created; status reaches READY & OK\\ \hline
Deploy an XGBoost model & Framework detected; booster materialized; served natively & OK\\ \hline
Predict with the wrong feature count & Readable \texttt{502} carrying the model server's message & OK\\ \hline
Delete a deployment & Cluster resources removed; record cleaned up & OK\\ \hline
\end{tabularx}
\end{table}

\begin{table}[H]\centering
\renewcommand{\arraystretch}{1.15}
\caption{Sprint 4 test scenarios --- Public API, monitoring, administration}\label{tab:tests-s4}
\small
\begin{tabularx}{\textwidth}{|X|X|c|}
\hline
\rowcolor{primary!15}
\textbf{Scenario} & \textbf{Expected result} & \textbf{Status}\\
\hline
Call the public endpoint with a valid key & Prediction returned without a platform account & OK\\ \hline
Call with a revoked key & \texttt{401} & OK\\ \hline
Send predictions, open the Monitoring tab & Volume, error rate, latency, per-origin and per-deployment breakdown reflect the calls & OK\\ \hline
Call an admin endpoint as a normal user & \texttt{403} & OK\\ \hline
Administrator attempts self-demotion & Refused with an explanatory error & OK\\ \hline
\end{tabularx}
\end{table}

% ---------------------------------------------------------------------------
\section{End-to-End Case Study: Telecom Churn}
\label{sec:case-study}

\subsection{Scenario}
The platform was validated with a real churn-prediction use case built on
operational data from a regional telecom operator: a file of terminated
customers and a file of network incidents, joined and transformed into
customer-level features by a Jupyter notebook. The complete journey runs
through the platform only (Figure~\ref{fig:casestudy-journey}): the notebook
and its two Excel files are uploaded as one zip; the notebook is
\textbf{automatically converted} to a script (Section~\ref{sec:sprint2});
the run executes on the cluster with autolog capturing every metric; the best
model is registered, promoted, deployed, and finally called from outside
through the public \gls{api}.

\begin{figure}[H]\centering
  \includegraphics[width=\textwidth,keepaspectratio]{casestudy_journey}
  \caption{Case-study journey: one dataset, the full lifecycle on the platform}
  \label{fig:casestudy-journey}
\end{figure}

\subsection{Evaluation Metrics}
Churn data is imbalanced: most clients stay, few leave. Three complementary
metrics are therefore used. \textbf{Accuracy} measures overall correctness but
can be inflated by the majority class. The \textbf{F1-score} focuses on the
churn class only, combining precision (are the raised alerts correct?) and
recall (how many real churners are caught?), so it cannot be gamed by
predicting ``no churn'' everywhere. The \textbf{\gls{roc}-\gls{auc}} evaluates
the model's \emph{ranking}: the probability that a random churner receives a
higher risk score than a random non-churner, independently of any decision
threshold --- the right measure when the business goal is ``call the most
at-risk clients first''.

\subsection{Model and Results}
The model is an XGBoost classifier trained on an 80/20 split, with SMOTE
applied to the training part to compensate for class imbalance.
Table~\ref{tab:churn-metrics} reports the evaluation results captured by the
platform, and Table~\ref{tab:churn-cm} the confusion matrix on the test set.

\begin{table}[H]\centering
\renewcommand{\arraystretch}{1.2}
\caption{Churn model: evaluation metrics (test set)}\label{tab:churn-metrics}
\small
\begin{tabularx}{\textwidth}{|X|c|X|}
\hline
\rowcolor{primary!15}
\textbf{Metric} & \textbf{Value} & \textbf{Reading}\\
\hline
Accuracy & 0.881 & Overall correctness on all clients\\ \hline
F1-score (churn class) & 0.773 & Balance of precision and recall on churners\\ \hline
ROC-AUC & 0.911 & Excellent ranking of clients by churn risk\\ \hline
ROC-AUC (5-fold CV) & $\approx$ 0.91 & Stable across folds --- no overfitting to one split\\ \hline
\end{tabularx}
\end{table}

\begin{table}[H]\centering
\renewcommand{\arraystretch}{1.2}
\caption{Churn model: confusion matrix (test set)}\label{tab:churn-cm}
\small
\begin{tabular}{|l|c|c|}
\hline
\rowcolor{primary!15}
 & \textbf{Predicted: active} & \textbf{Predicted: churn}\\
\hline
\textbf{Real: active} & 669 & 35\\ \hline
\textbf{Real: churn}  & 87  & 201\\ \hline
\end{tabular}
\end{table}

The business reading confirms the ranking quality measured by the
\gls{roc}-\gls{auc}: when clients are sorted by predicted risk, the low-risk
segment contains about 5\% real churners, the medium segment about 27\%, and
the high-risk segment about \textbf{96\%} --- exactly the list a retention
team should call first (Figure~\ref{fig:casestudy-proof}).

\begin{figure}[H]\centering
  \includegraphics[width=\textwidth,keepaspectratio]{casestudy_proof}
  \caption{Case-study results: real metrics and risk segments}
  \label{fig:casestudy-proof}
\end{figure}

\bigshot{screen_churn_deploy}{The churn model, registered and deployed}{The
champion run for this case study, registered from the leaderboard and deployed
as a live \texttt{InferenceService}: the platform's own registry and
deployment record for the exact model behind the metrics above, status
\texttt{READY}.}

\subsection{Consuming the Model from Outside}
To close the loop, the deployed churn model was called from an external
\gls{http} client using a per-deployment \gls{api} key, exactly as a
third-party retention application would (Listing~\ref{lst:churn-call}). The
15 values are the model's actual training features, in order ---
\texttt{nb\_reclamations}, \texttt{delai\_moy\_resolution},
\texttt{delai\_max\_resolution}, \texttt{delai\_min\_resolution},
\texttt{delai\_std\_resolution}, \texttt{nb\_incidents\_longs},
\texttt{nb\_non\_resolus}, \texttt{nb\_categories\_distinctes},
\texttt{a\_eu\_reclamation}, \texttt{ratio\_non\_resolu},
\texttt{ratio\_incidents\_longs}, \texttt{score\_risque}, \texttt{debit\_num},
\texttt{fsi\_encoded}, \texttt{cat\_encoded} --- and the response below is a
real capture, not an illustration. The call and its latency then appear in the
project's Monitoring tab.

\begin{lstlisting}[language=bash,caption={External prediction call against the deployed churn model (real request/response capture).},label={lst:churn-call}]
curl -X POST "https://mlops.firasmahjoubi.app/api/public/predict/<deployment_id>" \
  -H "Authorization: Bearer mlops_<key>" \
  -H "Content-Type: application/json" \
  -d '{"instances": [[5, 8.5, 15, 2, 4.2, 2, 1, 3, 1, 0.2, 0.4, 18.55, 8, 2, 1]]}'

{"predictions":[0.03778325021266937]}   % churn risk of this client: 3.8%
\end{lstlisting}

This exact call initially failed with a \texttt{502} carrying \emph{``training
data did not have the following fields: nb\_reclamations, \ldots''} ---
XGBoost's booster embeds the training \texttt{DataFrame}'s column names,
and the public/tester endpoints send a plain, unnamed array, so the model
server rejected every request for this deployment. The fix strips
\texttt{feature\_names}/\texttt{feature\_types} from the booster before it is
saved as the servable \texttt{model.bst} (Section~\ref{sec:sprint3}), so an
unnamed array matches what the booster now expects. Applied to this
deployment's artifact and the predictor pod restarted, the same call now
returns \texttt{200} with a real prediction, as captured above.

\bigshot{screen_churn_monitoring}{Live traffic from the external call}{The
churn deployment's Monitoring tab immediately after the Postman call above:
the request is counted, its latency measured, and it is attributed to the
public \gls{api} origin --- closing the loop from an outside \gls{http} client
to the platform's own serving telemetry.}

\subsection{What the Case Study Demonstrates}
This single use case exercises every layer built in the four sprints:
\begin{itemize}[leftmargin=1.4em]
  \item the \textbf{notebook conversion} handled a real, messy notebook
        (Excel dependencies, \texttt{!pip install}, French text) without
        manual editing;
  \item \textbf{autolog} captured all parameters, metrics, and evaluation
        artifacts with zero code changes;
  \item the \textbf{leaderboard} exposed the metrics above and flagged the
        champion run;
  \item \textbf{multi-framework serving} detected the XGBoost booster and
        materialized a servable file automatically;
  \item the deployed model answered predictions from the in-platform tester
        and from an external \gls{http} client through a hashed \gls{api}
        key, with every call visible in the \textbf{serving telemetry}.
\end{itemize}

% ---------------------------------------------------------------------------
\section{Conclusion}
The platform passed its functional scenarios across all four sprints, and a
real churn model traveled the full lifecycle --- upload, automatic conversion,
tracked training (\gls{roc}-\gls{auc} 0.911), registration, deployment, public
consumption, monitoring --- without a single line of infrastructure code
written by the data scientist. This is precisely the objective fixed in
Chapter~\ref{chap:context}.
